
This work introduces methodology for measuring how data curation strategies affect benchmark contamination. The Contamination Contribution Ratio (CCR) combines contamination detection with data attribution to quantify benchmark-specific influence from contaminated training examples. Methodology validation demonstrates:

\begin{itemize}
\item CCR scales linearly with contamination ($R^2$ = 0.9998)
\item Bootstrap statistics correctly detect CCR differences between filtering strategies
\item Removal intervention can establish causal importance of high-CCR examples
\item Influence Fragility Ratio distinguishes contaminated from non-contaminated examples
\end{itemize}

The primary limitation is that all experiments used simulated contamination due to GPU infrastructure constraints. Full empirical validation---training models on RedPajama-V2 with actual perplexity filtering and computing real TRAK attribution scores---is required to determine whether perplexity filtering naturally amplifies benchmark contamination in practice.

Future directions include:
\begin{itemize}
\item GPU cluster execution with actual model training
\item Span-level attribution analysis to investigate the weaker-than-expected IFR-redundancy correlation
\item Document-length stratification to rule out length as a confound
\item Verification of clean benchmark overlap assumptions
\item Scaling experiments to larger models
\end{itemize}
